% ============================================================
% f2014_beijing_science_q04_flowchart.tex
% 2014 北京卷（理科）第 4 题：程序框图
% 来源 PDF：普通高考/2014/2014北京理.pdf
% 原页：第 1 页（A3），印面 861，左栏第 4 题题干下方
% 结构：开始 → 输入 m, n 的值 → k=m, S=1 → 判断 k<m-n+1；
%       否则 S=S·k、k=k-1 后返回判断；是则输出 S，再结束。
% ============================================================
\begin{tikzpicture}[>=Stealth, line width=0.7pt, font=\small]
  % 开始
  \draw[rounded corners=6pt] (-0.95,0.45) rectangle (0.95,-0.45);
  \node at (0,0) {开始};
  \draw[->] (0,-0.45) -- (0,-1.05);

  % 输入
  \draw (-2.00,-1.05) -- (1.70,-1.05) -- (2.00,-2.05) -- (-1.70,-2.05) -- cycle;
  \node at (0,-1.55) {输入 $m$，$n$ 的值};
  \draw[->] (0,-2.05) -- (0,-2.55);

  % 初始化
  \draw (-1.55,-2.55) rectangle (1.55,-3.45);
  \node at (0,-3.00) {$k=m$，$S=1$};
  \draw[->] (0,-3.45) -- (0,-4.33);

  % 判断菱形
  \draw (0,-4.35) -- (2.20,-5.35) -- (0,-6.35) -- (-2.20,-5.35) -- cycle;
  \node at (0,-5.35) {$k<m-n+1$};

  % 是：输出
  \draw[->] (0,-6.35) -- (0,-6.85);
  \node at (-0.40,-6.60) {是};
  \draw (-2.00,-6.85) -- (1.60,-6.85) -- (2.00,-7.85) -- (-1.60,-7.85) -- cycle;
  \node at (0,-7.35) {输出 $S$};
  \draw[->] (0,-7.85) -- (0,-8.35);

  % 结束
  \draw[rounded corners=6pt] (-0.95,-8.35) rectangle (0.95,-9.25);
  \node at (0,-8.80) {结束};

  % 否：循环体
  \draw[->] (2.20,-5.35) -- (2.90,-5.35);
  \node at (2.52,-5.08) {否};
  \draw (2.90,-5.85) rectangle (5.10,-4.85);
  \node at (4.00,-5.35) {$S=S\cdot k$};
  \draw[->] (4.00,-4.85) -- (4.00,-4.15);
  \draw (2.90,-4.15) rectangle (5.10,-3.15);
  \node at (4.00,-3.65) {$k=k-1$};

  % 回边：由 k=k-1 向左汇入判断框上方的竖线
  \draw[->] (2.90,-3.65) -- (0,-3.65);
\end{tikzpicture}
